\documentclass[11pt]{article}
\usepackage[margin=1in]{geometry}
\usepackage{amsmath}
\usepackage{microtype}
\usepackage[hidelinks]{hyperref}

\newcommand{\Faq}{\textit{F.~aquilonia}}
\newcommand{\Fpo}{\textit{F.~polyctena}}
\newcommand{\Rst}{\ensuremath{R_{\mathrm{st}}}}

\title{\vspace{-3em}Among-population disequilibrium between unlinked ancestry
clusters tracks a shared ancestry axis, not pair-specific incompatibility}
\author{}
\date{}

\begin{document}
\maketitle
\vspace{-4em}

\noindent\textit{Draft results summary (Module D, the intrinsic arm). Methods are
deferred to Supplementary Materials. Findings are descriptive: the screen
characterises the observed distribution of among-population two-locus
disequilibrium and ranks candidate pairs, but a pair-specific incompatibility is
licensed only as an excess over the forthcoming recombination-matched neutral
null. All statements below concern that observed distribution.}

\section*{Results}

\paragraph{Design.}
An intrinsic incompatibility between two loci acts on their joint parental
ancestry, and across independent replicate populations should leave a pair-level
trace that per-locus sorting cannot show: the two loci fixing compatible ancestry
in the same direction repeatedly, so that their allele frequencies covary across
populations. Between loci that recombine freely this covariance cannot be produced
by linkage, and under drift alone independent replicates sort such loci
independently, so its neutral expectation is near zero. We define two clusters as
\emph{unlinked} on the genetic map---on different chromosomes, or on the same
chromosome more than $10$~cM apart---because admixture linkage disequilibrium
decays as $(1-c)^t$ irrespective of physical distance, and at the $\sim\!50$
generations of hybridisation implied for these populations $10$~cM corresponds to
$\sim\!99\%$ decay (and to roughly five mean ancestry-tract lengths). We screened
all such unlinked pairs of clusters that were polymorphic in the parents (pooled
parental MAF $\ge 0.15$; 27{,}223 clusters, $>\!3.6\times10^{8}$ pairs), measuring
the among-population covariance as \Rst{}, the correlation across the twenty
populations of a cluster's mean consensus dosage, and evaluating the exact Ohta
decomposition on the 401{,}945 pairs reaching $|\Rst| \ge 0.7$. Same-chromosome
pairs within $10$~cM (about $1.8$~million) are retained as a close-linkage
positive control.

\paragraph{Linkage raises among-population disequilibrium, and it decays to the
unlinked baseline by $\sim\!10$~cM.} The close-linked control behaves as it must:
same-chromosome pairs within $10$~cM carry a heavier tail of \Rst{} than unlinked
pairs, exceeding them by $1.8\times$, $2.9\times$ and $16\times$ at
$\Rst \ge 0.5$, $0.7$ and $0.9$. Resolving this by genetic distance shows the
expected decay: mean among-population $\Rst^2$ between same-chromosome clusters
falls from $0.082$ at $\sim\!1$~cM to the inter-chromosomal baseline of $0.058$ by
about $10$--$15$~cM ($0.061$ at $8$--$10$~cM). The empirical decorrelation scale
therefore matches the $\sim\!50$-generation expectation, and confirms from the data
that clusters more than $10$~cM apart carry no more among-population disequilibrium
than clusters on different chromosomes.

\paragraph{Among unlinked pairs, the decomposition points to structure, not
epistasis.} On the high-\Rst{} unlinked pairs the Ohta decomposition is unambiguous
about the source of the disequilibrium. The among-population component exceeds the
within-population component in virtually every pair (median $D^2_{ST} = 0.13$
versus $D^2_{IS} = 0.005$; $D^2_{ST} > D^2_{IS}$ in $>99.99\%$ of pairs), while the
disequilibrium of the averaged population is negligible (median
$D'^2_{ST} = 0.005$, below $D'^2_{IS}$ in every pair). A large among-population
variance component together with a vanishing average disequilibrium is the classic
signature of correlated allele-frequency \emph{divergence} among subpopulations---
drift and shared structure---rather than of a systematic two-locus association of
the kind epistatic selection would maintain within populations.

\paragraph{A single shared ancestry axis accounts for the pattern.}
The covariance is predominantly positive: $62\%$ of the high-\Rst{} unlinked pairs
covary in the same direction, and among pairs in which both clusters are themselves
directionally sorted, same-direction pairs outnumber opposite-direction pairs
$1.3$ to $1$. This is what a single genome-wide ancestry gradient produces---
populations that fixed \Faq{} ancestry did so across many diagnostic loci at once,
so those loci covary positively regardless of any interaction between them. The
among-population disequilibrium between unlinked clusters is thus dominated by the
same directional ancestry sorting characterised in the per-locus analyses, now seen
as the covariance it induces between loci, rather than by pair-specific coupling.

\paragraph{Interpretation and outlook.}
Before a neutral null, the screen finds no evidence of pair-specific incompatibility.
Linkage raises among-population disequilibrium and it decays away by $\sim\!10$~cM;
beyond that distance the disequilibrium that remains reads, by its own
decomposition, as correlated ancestry divergence along one genome-wide axis. This
is the expected descriptive outcome, and it sharpens rather than answers the
intrinsic question. A Dobzhansky--Muller interaction would appear as an
\emph{excess} of among-population disequilibrium, at specific unlinked pairs, over
what neutral sorting of the same replicate populations produces. Module D is
therefore run as a distribution-generator: it emits the observed \Rst{} and
among-population disequilibrium distributions, and candidate incompatibilities are
defined downstream as the pairs exceeding the recombination-matched neutral null,
evaluated by applying the identical pre-filter and decomposition to simulated
genotypes. That comparison, with the concurrent close-linkage reference, is the
decisive next step.

\end{document}
